\section{The serving stack and the workload it now carries}

The reorganization of the serving stack is not a proposal anyone needs to
make; it is happening. Mozilla.ai's articulation of an LLM control plane
names the three-plane split explicitly and maps it onto LLM
infrastructure \citep{mozilla2026}; NVIDIA's Dynamo separates scheduling and
placement from execution in a disaggregated runtime \citep{dynamo2025}; AIBrix
builds control-plane machinery as a first-class system layer
\citep{aibrix2025}; and external control planes with explicit admission control
are appearing for agentic workloads specifically \citep{mars2026}. What the new
layer provides, across these systems and the commercial gateways around
them, is a consistent portfolio: authentication and routing, rate limits,
per-request and per-session budget caps \citep{litellm2026,truefoundry2026},
priority scheduling with preemption \citep{vllm2023}, and reserved throughput
sold at the account level \citep{bedrock_ptu,azure_ptu}. Each mechanism is
useful. Together they constitute a control plane built on one implicit
assumption: that the unit of work is a request, that any request may be
delayed, preempted, migrated, or dropped under pressure, and that
fairness and efficiency are properties of request streams.

The workload violates the assumption. A flow, an agent executing a task,
a multi-step job with a declared objective, has five properties that
defeat specific mechanisms of the request-oriented portfolio. It spans
tens of invocations whose count is unknown at start, so per-invocation
latency optimization compounds queueing delay across the flow and tail
completion is governed by the worst episode. It accumulates key-value
cache and tool state that are its progress, so preemption with
recompute-on-resume or swap-and-requeue does not delay the work but
destroys it. It idles between invocations for seconds to minutes while
tools and humans respond, and under request accounting an idle flow's
resources are fair game, so resumption becomes randomized restart; the
idle gap is also, we note at the outset, the strongest argument against
the discipline this paper proposes and prices. Its value is concentrated
at completion, so a flow abandoned at ninety-five percent of its
invocations delivers nothing at ninety-five percent of its cost,
inverting the economics that make preemption cheap for divisible work.
And its consumption is unbounded by default, an agent in a defective or
adversarial loop drawing tokens across invocations that per-request caps
structurally cannot see, while session-level budget ceilings
\citep{litellm2026,truefoundry2026} answer the accumulation with rejection,
converting the entire prior spend into waste at the moment of
enforcement.

The portfolio, taken together, provides caps that see flows but cannot
feed them, schedulers that feed work but cannot promise it anything, and
account-level floors that promise capacity but cannot see flows. Nothing
guarantees that a specific task, admitted deliberately, will finish or be
deliberately stopped; and because no mechanism binds authority to a
specific flow, the portfolio offers no enforcement point for governing
one. The missing object is a completion-oriented resource contract:
resources sufficient to finish, granted at admission and not silently
withdrawn; a bound enforced as part of the same contract rather than as
an ambush at a ceiling; and a designated authority, outside the flow,
entitled to stop it. Section 4 specifies a discipline providing exactly
these clauses. Whether the contract is worth its price is the question
the rest of the paper answers.
